
\documentclass{article}
\usepackage{colortbl}
\usepackage{makecell}
\usepackage{multirow}
\usepackage{supertabular}

\begin{document}

\newcounter{utterance}

\centering \large Interaction Transcript for game `cladder', experiment `full\_v1.5\_default', episode 2103 with selfrefine{-}Qwen3.5{-}9B.
\vspace{24pt}

{ \footnotesize  \setcounter{utterance}{1}
\setlength{\tabcolsep}{0pt}
\begin{supertabular}{c@{$\;$}|p{.15\linewidth}@{}p{.15\linewidth}p{.15\linewidth}p{.15\linewidth}p{.15\linewidth}p{.15\linewidth}}
    \# & \multicolumn{2}{c}{Player} && \multicolumn{2}{c}{Game Master} \\
    \hline

    \theutterance \stepcounter{utterance}  
    & & & \multicolumn{4}{p{0.6\linewidth}}{
        \cellcolor[rgb]{0.9,0.9,0.9}{
            \makecell[{{p{\linewidth}}}]{
                \texttt{\tiny{[P1$\langle$GM]}}
                \texttt{You are an expert in causal inference. The following question is not a typical commonsense query, but rather a meticulously designed question created by a professor specializing in causal inference, intended to assess the students' mastery of the course content.} \\
\texttt{Imagine a self{-}contained, hypothetical world with only the following conditions, and without any unmentioned factors or causal relationships: Gender has a direct effect on department competitiveness and admission status. Department competitiveness has a direct effect on admission status. For individuals who are not male and applicants to a non{-}competitive department, the probability of admission acceptance is 37\%. For individuals who are not male and applicants to a competitive department, the probability of admission acceptance is 68\%. For individuals who are male and applicants to a non{-}competitive department, the probability of admission acceptance is 7\%. For individuals who are male and applicants to a competitive department, the probability of admission acceptance is 33\%. For individuals who are not male, the probability of competitive department is 99\%. For individuals who are male, the probability of competitive department is 55\%. Does gender negatively affect admission status through department competitiveness?} \\
\texttt{Start your answer with 'Yes' or 'No', followed by additional reasoning or evidence to support your explanation.} \\
            }
        }
    }
     \\ \\

    \theutterance \stepcounter{utterance}  
    & \multicolumn{4}{p{0.6\linewidth}}{
        \cellcolor[rgb]{0.561,0.737,0.561}{
            \makecell[{{p{\linewidth}}}]{
                \texttt{\tiny{[P1$\rangle$GM]}}
                \texttt{No} \\
\\ 
\texttt{**Reasoning:**} \\
\\ 
\texttt{To determine if gender negatively affects admission status *through* department competitiveness, we must calculate the **Natural Indirect Effect (NIE)** of changing gender from Male to Female, while holding department competitiveness constant.} \\
\\ 
\texttt{Let \$G\$ denote Gender (Male = 1, Female = 0) and \$D\$ denote Department Competitiveness (Competitive = 1, Non{-}competitive = 0). Let \$Y\$ denote Admission Status.} \\
\texttt{The question asks if the indirect effect of changing \$G\$ from 1 to 0 (Male \$\textbackslash{}to\$ Female) on \$Y\$ via \$D\$ is negative.} \\
\texttt{Mathematically, we compare \$E{[}Y | G=0, D=d{]}\$ vs \$E{[}Y | G=1, D=d{]}\$ for each level of \$d\$.} \\
\\ 
\texttt{**Step 1: Extract the Conditional Probabilities (\$P(Y|G, D)\$)**} \\
\texttt{From the prompt:} \\
\texttt{*   \$P(Y=1 | G=0, D=0) = 37\textbackslash{}\%\$ (Female, Non{-}competitive)} \\
\texttt{*   \$P(Y=1 | G=0, D=1) = 68\textbackslash{}\%\$ (Female, Competitive)} \\
\texttt{*   \$P(Y=1 | G=1, D=0) = 7\textbackslash{}\%\$  (Male, Non{-}competitive)} \\
\texttt{*   \$P(Y=1 | G=1, D=1) = 33\textbackslash{}\%\$  (Male, Competitive)} \\
\\ 
\texttt{**Step 2: Analyze the Indirect Effect for each level of Department Competitiveness**} \\
\texttt{The indirect effect measures the change in admission probability when gender changes from Male to Female, *holding the department competitiveness fixed*.} \\
\\ 
\texttt{*   **Scenario A: Holding Department as Non{-}competitive (\$D=0\$)**} \\
\texttt{    *   Probability for Female: 37\%} \\
\texttt{    *   Probability for Male: 7\%} \\
\texttt{    *   Effect of changing Male \$\textbackslash{}to\$ Female: \$37\textbackslash{}\% {-} 7\textbackslash{}\% = +30\textbackslash{}\%\$} \\
\texttt{    *   *Result:* The effect is **positive**. Being female increases admission chances in non{-}competitive departments.} \\
\\ 
\texttt{*   **Scenario B: Holding Department as Competitive (\$D=1\$)**} \\
\texttt{    *   Probability for Female: 68\%} \\
\texttt{    *   Probability for Male: 33\%} \\
\texttt{    *   Effect of changing Male \$\textbackslash{}to\$ Female: \$68\textbackslash{}\% {-} 33\textbackslash{}\% = +35\textbackslash{}\%\$} \\
\texttt{    *   *Result:* The effect is **positive**. Being female increases admission chances in competitive departments.} \\
\\ 
\texttt{**Conclusion:**} \\
\texttt{In both possible states of the mediator (department competitiveness), changing gender from Male to Female results in a **higher** probability of admission. Since the indirect effect is positive in all scenarios, gender does **not**} \\
            }
        }
    }
    & & \\ \\

    \theutterance \stepcounter{utterance}  
    & & & \multicolumn{2}{p{0.3\linewidth}}{
        \cellcolor[rgb]{0.9,0.9,0.9}{
            \makecell[{{p{\linewidth}}}]{
                \texttt{\tiny{[GM$|$GM]}}
                \texttt{No} \\
            }
        }
    }
    & & \\ \\

    \theutterance \stepcounter{utterance}  
    & & & \multicolumn{2}{p{0.3\linewidth}}{
        \cellcolor[rgb]{0.9,0.9,0.9}{
            \makecell[{{p{\linewidth}}}]{
                \texttt{\tiny{[GM$|$GM]}}
                \texttt{yes} \\
            }
        }
    }
    & & \\ \\

    \theutterance \stepcounter{utterance}  
    & & & \multicolumn{2}{p{0.3\linewidth}}{
        \cellcolor[rgb]{0.9,0.9,0.9}{
            \makecell[{{p{\linewidth}}}]{
                \texttt{\tiny{[GM$|$GM]}}
                \texttt{game\_result = LOSE} \\
            }
        }
    }
    & & \\ \\

\end{supertabular}
}

\end{document}
